\chapter{Coherence of locally trivial line bundles}
\label{chap:Picard_LineBundleCoherence}
% archon:covers AlgebraicJacobian/Picard/LineBundleCoherence.lean

% SCOPE. This chapter proves that a locally trivial line bundle is coherent:
% \(\texttt{IsLocallyTrivial} M \Rightarrow M.\mathtt{IsFinitePresentation}\),
% together with a rank-one / flatness record. The general implication
% ``\(M\) tensor-invertible \(\Rightarrow\) \(M\) locally free of rank one''
% (the converse spreading-out, Stacks Lemma~0B8M when stalks are local) is
% \emph{not} crossed here: the carrier is already locally trivial, so the
% finite-presentation spreading-out is never needed.

\section{Setup and motivation}
\label{sec:lbc_setup}

Fix a scheme \(X\) and write \(X.\mathtt{Modules}\) for the category of sheaves
of \(\mathcal{O}_X\)-modules (Mathlib's \(\mathtt{SheafOfModules}\) on the
ringed-space site of \(X\), with structure sheaf \(X.\mathtt{ringCatSheaf}\)).
A \emph{line bundle} on \(X\) is an invertible \(\mathcal{O}_X\)-module; in the
sense of \cref{chap:Picard_LineBundlePullback} we model line bundles by the
predicate \(\texttt{IsLocallyTrivial}\): a module \(M : X.\mathtt{Modules}\) is locally
trivial of rank one if every point \(x \in X\) has an \emph{affine} open
neighbourhood \(U\) on which the restriction \(M|_U\) is isomorphic, as an
\(\mathcal{O}_U\)-module, to the structure sheaf-as-module \(\struct{U}\):
\[
  \texttt{IsLocallyTrivial}(M) \;:=\;
  \forall x \in X,\ \exists\, U \ni x \text{ affine open},\quad
  M|_U \;\cong\; \struct{U}.
\]
This is the rank-one case of Stacks Definition~17.25.1
(Tag~01CS, invertible \(\mathcal{O}_X\)-modules), and
agrees with the classical ``locally free of rank one'' description by Stacks
Lemma~17.25.4 (Tag~0B8M).

The goal of the chapter is to record that such a module is \emph{coherent},
namely finitely presented:
\(M.\mathtt{IsFinitePresentation}\) in Mathlib's
\(\mathtt{SheafOfModules.IsFinitePresentation}\) (built from a
\(\mathtt{QuasicoherentData}\)). The mathematics is elementary: on each
trivialising chart \(M|_U \cong \struct{U}\) is finite free of rank one, hence
has the trivial finite presentation, and finite presentation is by definition a
local condition. No stalk computation and no spreading-out are involved. The
content is the bookkeeping that turns the pointwise trivialisation into the
indexed \(\mathtt{QuasicoherentData}\) Mathlib expects.

\section{Site instances for quasi-coherent data}
\label{sec:lbc_site_instances}

% NOTE (chief unknown). Mathlib's SheafOfModules.IsFinitePresentation /
% QuasicoherentData are stated under three site hypotheses on the slice
% topologies J.over X of the underlying site J of X.ringCatSheaf:
%   [∀ x, (J.over x).HasSheafCompose (forget₂ RingCat AddCommGrpCat)]
%   [∀ x, HasWeakSheafify (J.over x) AddCommGrpCat]
%   [∀ x, (J.over x).WEqualsLocallyBijective AddCommGrpCat]
% These must be discharged (or assumed as instances) before the statements in
% this chapter typecheck for X.ringCatSheaf. The project already lives in the
% SheafOfModules X.ringCatSheaf world, so these should resolve; this is the one
% item to verify when the Lean file is scaffolded.

The Mathlib predicates \(\mathtt{SheafOfModules.IsFinitePresentation}\) and
\(\mathtt{SheafOfModules.QuasicoherentData}\) are stated relative to a
Grothendieck topology \(J\) on a base category \(C\) and a sheaf of rings
\(R : \Sheaf(J, \RingCat)\). They require, for every object \(x : C\), three
instances on the slice topology \(J/x\) (Mathlib \(J.\mathtt{over}\ x\)):
\[
  (J/x).\mathtt{HasSheafCompose}(\forget_{\RingCat \to \AddCommGrpCat}),\quad
  \HasWeakSheafify(J/x, \AddCommGrpCat),\quad
  (J/x).\mathtt{WEqualsLocallyBijective}(\AddCommGrpCat).
\]
For \(X.\mathtt{Modules} = \mathtt{SheafOfModules}\ X.\mathtt{ringCatSheaf}\)
these are the ambient site instances of the ringed space underlying \(X\). We
record their availability as a standing assumption for the chapter; verifying it
for \(X.\mathtt{ringCatSheaf}\) is the single instance-resolution check the
formalisation must clear at entry, and is independent of the mathematical
content below.

\begin{remark}
  \label{rem:lbc_site_instances}
  \dcref{custom}
  For the ringed-space site underlying \(X\) with sheaf of rings
  \(R = X.\mathtt{ringCatSheaf}\), every slice topology \(J/x\) carries the
  three instances \(\mathtt{HasSheafCompose}\), \(\HasWeakSheafify\) and
  \(\mathtt{WEqualsLocallyBijective}\) (for the forgetful functor
  \(\RingCat \to \AddCommGrpCat\)) demanded by
  \(\mathtt{SheafOfModules.QuasicoherentData}\). Under this assumption the
  predicate \(M.\mathtt{IsFinitePresentation}\) is well-formed for every
  \(M : X.\mathtt{Modules}\).
\end{remark}

\section{From pointwise triviality to a trivialising cover}
\label{sec:lbc_cover}

\begin{lemma}[Trivialising affine cover]\leanok
  \label{lem:lbc_trivializing_cover}
  \dcref{custom}
  \lean{AlgebraicGeometry.Scheme.LineBundle.IsLocallyTrivial.exists_trivializing_cover}
  \uses{def:line_bundle_on_product}
  % SOURCE: \cite{stacks-project}, Modules on Ringed Spaces, Definition 17.25.1
  % "invertible module" (Tag 01CS), Section 17.25 (Tag 01CR)
  % (read from references/stacks-modules.tex, L4046-L4059).
  % SOURCE QUOTE:
  % "Let $(X, \mathcal{O}_X)$ be a ringed space. An
  %  {\it invertible $\mathcal{O}_X$-module} is a sheaf
  %  of $\mathcal{O}_X$-modules $\mathcal{L}$ such that
  %  the functor
  %  $$
  %  \textit{Mod}(\mathcal{O}_X) \longrightarrow \textit{Mod}(\mathcal{O}_X),\quad
  %  \mathcal{F} \longmapsto \mathcal{L} \otimes_{\mathcal{O}_X} \mathcal{F}
  %  $$
  %  is an equivalence of categories. We say that $\mathcal{L}$ is
  %  {\it trivial} if it is isomorphic as an $\mathcal{O}_X$-module
  %  to $\mathcal{O}_X$."
  \textit{Source: \cite{stacks-project}, Tag~01CS (Modules, Definition~17.25.1).}
  Let \(M : X.\mathtt{Modules}\) be locally trivial, \(\texttt{IsLocallyTrivial}(M)\).
  Then there exist an index set \(I\), a family of affine open subschemes
  \(\{U_i\}_{i \in I}\) of \(X\) covering \(X\) (i.e. \(\bigcup_i U_i = X\)), and
  for each \(i\) an isomorphism of \(\mathcal{O}_{U_i}\)-modules
  \[
    e_i \;:\; M|_{U_i} \;\xrightarrow{\ \sim\ }\; \struct{U_i}
  \]
  to the structure sheaf-as-module on \(U_i\).
\end{lemma}

\begin{proof}
  \uses{def:line_bundle_on_product}
  \leanok
  Apply the hypothesis \(\texttt{IsLocallyTrivial}(M)\) at every point \(x \in X\): it
  produces an affine open \(U_x \ni x\) and an isomorphism
  \(e_x : M|_{U_x} \cong \struct{U_x}\). Take the index set \(I := X\) (the
  underlying set of points), and to the index \(x\) assign the chosen
  neighbourhood \(U_x\) and the chosen isomorphism \(e_x\). The family
  \(\{U_x\}_{x \in X}\) covers \(X\) because each \(x\) lies in its own \(U_x\).
  This is exactly the data the definition supplies, packaged as an indexed
  family rather than a pointwise existential.
\end{proof}

\section{The canonical presentation of the unit module}
\label{sec:lbc_unit_presentation}

The chart presentation of \cref{sec:lbc_chart_presentation} is obtained by
transporting a fixed canonical presentation of the unit module
\(\struct{} = \mathtt{SheafOfModules.unit}\,R\) along the trivialising
isomorphism. We record that canonical presentation here. The unit module is free
of rank one, so it has a presentation with a single generator and no relations.

\begin{definition}[Free–unit isomorphism]\leanok
  \label{def:lbc_free_unit_iso}
  \dcref{custom}
  \lean{AlgebraicGeometry.Scheme.LineBundle.freeUnitIso}
  Over a sheaf of rings \(R\) on a site, the free \(R\)-module on a single
  generator is canonically isomorphic to the unit module:
  \[
    \mathtt{free}\,(\mathrm{PUnit}) \;\xrightarrow{\ \sim\ }\; \mathtt{unit}\,R,
  \]
  the coproduct of a one-element family of copies of \(\mathtt{unit}\,R\) being
  \(\mathtt{unit}\,R\) itself.
\end{definition}

\begin{proof}

\leanok
  Proved directly in Lean.
\end{proof}

\begin{definition}[Unit generating section]\leanok
  \label{def:lbc_unit_generators}
  \dcref{custom}
  \lean{AlgebraicGeometry.Scheme.LineBundle.unitGenerators}
  \uses{def:lbc_free_unit_iso}
  The canonical generating family of the unit module \(\mathtt{unit}\,R\): the
  single section indexed by \(\mathrm{PUnit}\), namely the image of the free
  generator under the free–unit isomorphism \cref{def:lbc_free_unit_iso}. The
  associated map \(\mathtt{unit}\,R^{\oplus 1} \to \mathtt{unit}\,R\) is an
  epimorphism, so the family generates.
\end{definition}

\begin{proof}

\leanok
  Proved directly in Lean.
\end{proof}

\begin{definition}[Canonical unit presentation]\leanok
  \label{def:lbc_unit_presentation}
  \dcref{custom}
  \lean{AlgebraicGeometry.Scheme.LineBundle.unitPresentation}
  \uses{def:lbc_unit_generators}
  The canonical finite free presentation of the unit module \(\mathtt{unit}\,R\):
  one generator (\cref{def:lbc_unit_generators}) and no relations (indexed by the
  empty type). Since the generating map \(\mathtt{unit}\,R^{\oplus 1} \to
  \mathtt{unit}\,R\) is an isomorphism, its relation kernel is the zero object,
  so the empty family of relations suffices. This presentation is finite.
\end{definition}

\begin{proof}

\leanok
  Proved directly in Lean.
\end{proof}

\section{The trivial presentation on a chart}
\label{sec:lbc_chart_presentation}

\begin{lemma}[Finite free presentation on a chart]\leanok
  \label{lem:lbc_chart_presentation}
  \dcref{custom}
  \lean{AlgebraicGeometry.Scheme.LineBundle.IsLocallyTrivial.chartPresentation}
  \uses{lem:lbc_trivializing_cover, def:lbc_unit_presentation,
    def:linebundle_chart_over_iso}
  % SOURCE: \cite{stacks-project}, Modules on Ringed Spaces, section
  % "Modules of finite presentation" (\label{definition-finite-presentation})
  % (read from references/stacks-modules.tex, L1377-L1392).
  % SOURCE QUOTE:
  % "We say that $\mathcal{F}$ is of {\it finite presentation}
  %  if for every point $x \in X$ there exists an open neighbourhood
  %  $x\in U \subset X$, and  $n, m \in \mathbf{N}$ such that $\mathcal{F}|_U$
  %  is isomorphic to the cokernel of a map
  %  $$
  %  \bigoplus\nolimits_{j = 1, \ldots, m}
  %  \mathcal{O}_U
  %  \longrightarrow
  %  \bigoplus\nolimits_{i = 1, \ldots, n}
  %  \mathcal{O}_U
  %  $$"
  \textit{Source: \cite{stacks-project}, Modules, ``Modules of finite
  presentation''.}
  On each chart \(U_i\) of the trivialising cover of
  \cref{lem:lbc_trivializing_cover}, the restriction \(M|_{U_i}\) admits the
  trivial finite free presentation
  \[
    0 \;\longrightarrow\; \struct{U_i} \;\longrightarrow\; M|_{U_i}
    \;\longrightarrow\; 0,
  \]
  i.e. \(M|_{U_i}\) is the cokernel of the zero map
  \(\bigoplus_{j \in \varnothing} \struct{U_i} \to \struct{U_i}\)
  (the case \(m = 0\), \(n = 1\) of the finite-presentation pattern). The
  unit sheaf-as-module \(\struct{U_i}\) (Mathlib's \(\mathtt{SheafOfModules.unit}\))
  carries a canonical free presentation on one generator with no relations;
  this presentation is finite.
\end{lemma}

\begin{proof}
  \uses{lem:lbc_trivializing_cover}
  \leanok
  The structure sheaf-as-module \(\struct{U_i} = \mathtt{SheafOfModules.unit}\,
  (U_i).\mathtt{ringCatSheaf}\) is free of rank one on the single generator
  \(1\): the identity map \(\struct{U_i} \to \struct{U_i}\) exhibits it as
  \(\bigoplus_{i \in \{*\}} \struct{U_i}\), so the surjection
  \(\struct{U_i}^{\oplus 1} \twoheadrightarrow \struct{U_i}\) is an isomorphism
  and its kernel is zero, generated by the empty family. Thus \(\struct{U_i}\)
  has the canonical presentation with \(n = 1\) generators and \(m = 0\)
  relations; this canonical \(\mathtt{unit}\)-presentation is the rank-one free
  presentation, and Mathlib records it as \emph{finite} (the
  \(\mathtt{SheafOfModules.unit}\) finite-presentation instance).

  The chart presentation is obtained by transporting the canonical finite
  presentation of \(\mathtt{unit}\,(U_i).\mathtt{ringCatSheaf}\) along the
  isomorphism \(e_i\), using \(\mathtt{SheafOfModules.Presentation.ofIsIso}\).
  This construction carries any finite presentation of a module to a finite
  presentation of any isomorphic module, so the transported presentation of
  \(M|_{U_i}\) is again finite.
  % NOTE. The over-restrict categorical bridge is the shared construction
  % `Scheme.Modules.chartOverIso` (lem:chart_over_iso,
  % chap:Picard_SheafOverEquivalence).
\end{proof}

\section{Finite presentation of a locally trivial line bundle}
\label{sec:lbc_finite_presentation}

\begin{theorem}[Locally trivial \(\Rightarrow\) finitely presented]\leanok
  \label{thm:lbc_isFinitePresentation}
  \dcref{custom}
  \lean{AlgebraicGeometry.Scheme.LineBundle.IsLocallyTrivial.isFinitePresentation}
  \uses{lem:lbc_trivializing_cover, lem:lbc_chart_presentation, rem:lbc_site_instances}
  % SOURCE: \cite{stacks-project}, Modules on Ringed Spaces,
  % Lemma 17.25.4 (Tag 0B8M), invertible vs locally free of rank one
  % (read from references/stacks-modules.tex, L4159-L4165).
  % SOURCE QUOTE:
  % "Let $(X, \mathcal{O}_X)$ be a ringed space. Any locally free
  %  $\mathcal{O}_X$-module of rank $1$ is invertible.
  %  If all stalks $\mathcal{O}_{X, x}$ are local rings, then
  %  the converse holds as well (but in general this is not the case)."
  \textit{Source: \cite{stacks-project}, Tag~0B8M (Modules, Lemma~17.25.4).}
  Let \(X\) be a scheme satisfying \cref{rem:lbc_site_instances} and let
  \(M : X.\mathtt{Modules}\) be locally trivial, \(\texttt{IsLocallyTrivial}(M)\). Then
  \(M\) is finitely presented: \(M.\mathtt{IsFinitePresentation}\) holds in
  Mathlib's \(\mathtt{SheafOfModules.IsFinitePresentation}\).
\end{theorem}

\begin{proof}
  \uses{lem:lbc_trivializing_cover, lem:lbc_chart_presentation}
  \leanok
  By \cref{lem:lbc_trivializing_cover} choose a trivialising affine cover
  \(\{U_i\}_{i \in I}\) with isomorphisms \(e_i : M|_{U_i} \cong \struct{U_i}\).
  Assemble these into a \(\mathtt{QuasicoherentData}\) for \(M\): the index type
  is \(I\), the covering objects are the affine opens \(U_i\) (which cover \(X\)),
  and the presentation attached to index \(i\) is the finite free presentation of
  \(M|_{U_i}\) (equivalently, of the restriction \(M.\mathtt{over}\,U_i\))
  produced in \cref{lem:lbc_chart_presentation}. Each such presentation is
  finite by \cref{lem:lbc_chart_presentation}; the assembled
  \(\mathtt{QuasicoherentData}\) therefore satisfies
  \(\mathtt{QuasicoherentData.IsFinitePresentation}\), and the constructor
  \(\mathtt{SheafOfModules.IsFinitePresentation.mk}\) yields
  \(M.\mathtt{IsFinitePresentation}\).

  The only non-formal input is the trivialising cover; everything else is the
  definitional unwinding of finite presentation as a local property
  (Stacks, ``Modules of finite presentation''), since a rank-one free module is its own finite
  presentation. No stalk computation and no neighbourhood-spreading argument is
  used: the hypothesis already provides the local free models.
\end{proof}

\begin{corollary}[Locally trivial \(\Rightarrow\) finite type]\leanok
  \label{cor:lbc_isFiniteType}
  \dcref{custom}
  \lean{AlgebraicGeometry.Scheme.LineBundle.IsLocallyTrivial.isFiniteType}
  \uses{thm:lbc_isFinitePresentation}
  Under the hypotheses of \cref{thm:lbc_isFinitePresentation}, the locally
  trivial module \(M\) is of finite type (\(M.\mathtt{IsFiniteType}\)). Quasi-coherence
  (\(M.\mathtt{IsQuasicoherent}\)) follows from finite presentation via
  \(\mathtt{SheafOfModules.instIsQuasicoherentOfIsFinitePresentation}\).
\end{corollary}

\begin{proof}
  \uses{thm:lbc_isFinitePresentation}
  \leanok
  Finite presentation implies finite type by Mathlib's
  \(\mathtt{SheafOfModules.instIsFiniteTypeOfIsFinitePresentation}\), immediate
  from \cref{thm:lbc_isFinitePresentation}. Quasi-coherence is immediate from
  finite presentation by
  \(\mathtt{SheafOfModules.instIsQuasicoherentOfIsFinitePresentation}\).
\end{proof}

\section{Rank one and flatness record}
\label{sec:lbc_rank_flat}

% NOTE. Rank-one freeness and flatness are local conditions, so they are
% recorded chart-locally, on the trivialising charts themselves, rather than
% through a global predicate.

\begin{lemma}[Chart-local rank one and flatness]\leanok
  \label{lem:lbc_rank_flat}
  \dcref{custom}
  \lean{AlgebraicGeometry.Scheme.LineBundle.IsLocallyTrivial.chart_free_rank_one}
  \uses{lem:lbc_trivializing_cover}
  % SOURCE: \cite{stacks-project}, Modules on Ringed Spaces, section
  % "Locally free sheaves" (\label{definition-locally-free}, "(finite)
  % locally free of rank $r$")
  % (read from references/stacks-modules.tex, L2079-L2094).
  % SOURCE QUOTE:
  % "\item We say $\mathcal{F}$ is {\it locally free} if for every
  %  point $x \in X$ there exist a set $I$ and an open
  %  neighbourhood $x \in U \subset X$
  %  such that $\mathcal{F}|_U$ is isomorphic to
  %  $\bigoplus_{i \in I} \mathcal{O}_X|_U$ as an $\mathcal{O}_X|_U$-module.
  %  \item We say $\mathcal{F}$ is {\it finite locally free} if we may
  %  choose the index sets $I$ to be finite.
  %  \item We say $\mathcal{F}$ is {\it finite locally free of rank $r$}
  %  if we may choose the index sets $I$ to have cardinality $r$."
  \textit{Source: \cite{stacks-project}, Modules, ``Locally free sheaves'' (finite
  locally free of rank \(r\)).}
  Let \(M : X.\mathtt{Modules}\) be locally trivial. The lemma records, at each
  point \(x \in X\), a chart \(U \ni x\) together with the trivialising iso
  \(M|_U \cong \struct{U}\) (to \(\mathtt{SheafOfModules.unit}\)). This iso is the
  rank-one free \emph{model}: rank-one freeness and flatness are
  \emph{consequences} of this iso, because \(\struct{U}\) is the free
  \(\mathcal{O}_U\)-module of rank one and a flat module over itself (a ring is
  flat over itself), and both properties transport along \(e\). In the sense of
  \cref{def:line_bundle_on_product} this exhibits \(M\) as finite locally free
  of rank one and flat on the cover.
\end{lemma}

\begin{proof}
  \uses{lem:lbc_trivializing_cover}
  \leanok
  Immediate from the trivialisation: on \(U_i\), \(e_i\) identifies \(M|_{U_i}\)
  with \(\struct{U_i} = \bigoplus_{i \in \{*\}} \struct{U_i}\), the free module
  of rank one. The structure sheaf-as-module is flat over itself, and an
  isomorphic module is flat; the index set \(\{*\}\) has cardinality one, so the
  rank is one. The properties are recorded chart-locally because rank-one
  freeness and flatness are local conditions, so the chart-local data determine
  them.
\end{proof}
